\chapter{The generative setting}\label{ch:grammars}

This chapter fixes the objects of the monograph: generative systems and their histories, the UNNS operator SELECT-II as the principal realisation, symbolic event grammars as a second realisation, and the permutohedral complexes that carry the topology. It ends with a contract stating which constructions need integer divisibility and which hold for arbitrary finite sets of histories.

\section{Generative systems, histories and fibres}
\begin{definition}[generative system]\label[definition]{def:grammar}
A \emph{generative system} $\mathcal G=(\Lambda,X,X_0,(\tau_\ell)_{\ell\in\Lambda})$ consists of a set $\Lambda$ of \emph{labels}, a set $X$ of \emph{states}, a set $X_0\subseteq X$ of \emph{sources}, and for each label a partial map $\tau_\ell\colon X\rightharpoonup X$. A \emph{history} is a pair $h=(x_0;\ell_1,\dots,\ell_n)$ with $x_0\in X_0$ such that $x_i=\tau_{\ell_i}(x_{i-1})$ is defined for $1\le i\le n$; the word $w(h)=(\ell_1,\dots,\ell_n)$ is its \emph{word}, $x_n$ its \emph{endpoint} $t(h)$, and the multiset $\{\ell_1,\dots,\ell_n\}$ its \emph{content} $Q(h)$. The \emph{fibre} of $t\in X$ is the set $F_t$ of histories with endpoint $t$, and $F_t(x_0)$ the subset with source $x_0$. When $F_t$ is finite and nonempty, the \emph{(Hartley) ambiguity} of $t$ is $A(t)=\log_2|F_t|$.
\end{definition}

A fibre is a disjoint union of \emph{pieces} $P(x_0,Q)$ of histories with a given source and content. We write $\Ent(\cdot)$ for Shannon entropy in bits, and use the uniform distribution on a finite fibre unless another is stated.

\begin{definition}[commutative realization, legal swap]\label[definition]{def:swap}
A generative system is a \emph{commutative realization} if $X$ lies in an abelian group (written multiplicatively) and $\tau_\ell(x)=x\,r_\ell$ for a fixed $r_\ell$ on a domain $\mathrm{dom}\,\tau_\ell\subseteq X$. Then $t(h)=x_0\prod_i r_{\ell_i}$ depends only on the source and content. A \emph{legal adjacent swap} exchanges two adjacent distinct labels $\ell_i\ne\ell_{i+1}$ of a history when the exchanged word is still a history from the same source. The \emph{swap graph} $G_t$ of a fibre has the histories as vertices and the legal swaps as edges; its connected components are the \emph{recombination classes}, and the class of $h$ is written $C(h)$.
\end{definition}

\begin{lemma}\label[lemma]{lem:swap-invariants}
In a commutative realization a legal swap preserves source, content and endpoint. Hence the random variables $T$ (endpoint), $(M,Q)$ (source and content), $C$ (class) and $W$ (the history) satisfy that each determines the previous one: $T=f_1(M,Q)$, $(M,Q)=f_2(C)$, $C=f_3(W)$.
\end{lemma}
\begin{proof}
Swapping two labels permutes the word, so the multiset of labels and the source are unchanged, and $t=x_0\prod r_{\ell_i}$ is unchanged by commutativity of the group. Source and content are therefore constant on connected components of $G_t$, hence functions of the class; the class is a function of the history; the endpoint is a function of source and content.
\end{proof}

\section{The UNNS operator SELECT-II}\label{sec:selectII}
SELECT-II is the operator of the Substrate whose reachability and complete-history structure was established in the Foundations of the Substrate~\cite{UNNSFound}. Its histories can be enumerated exactly and its admissibility is governed by divisibility, which makes it a rich source of finite combinatorial structure. It is the principal worked realisation of this monograph; the monograph does not claim that it is the only generative rule of the Substrate.

\begin{definition}[SELECT-II]\label[definition]{def:selectII}
States are positive integers, labels are positive integers $N\ge1$, and
\begin{equation}\label[definition]{eq:selectII}
T_N(m)=m\,R(N),\qquad R(N)=\frac{(N+1)^2}{N}=N+2+\frac1N,\qquad \mathrm{dom}\,T_N=\{m: N\mid m\}.
\end{equation}
A source is a positive integer. A history is \emph{complete} if its source is squarefree (we show in \cref{prop:fibre-finite} that this means the source has no predecessor). The \emph{complete-history fibre} $F_t$ consists of all complete histories ending at $t$; the \emph{fixed-source fibre} $F_t(m_0)$ of all histories from $m_0$.
\end{definition}
SELECT-II is a commutative realization with $r_N=R(N)$, so \cref{lem:swap-invariants} applies. A label word is \emph{formally} applicable from any seed, since $R(N)$ is rational, but \emph{admissible} only if every intermediate state is an integer divisible by the next label.

\section{Symbolic event grammars}\label{sec:symbolic}
For the topological chapters no integer state is needed. Take $k$ distinguishable events and a set $A\subseteq S_k$ of admissible orderings. Two families occur. A \emph{pairwise grammar} is given by a digraph $\Gamma$ on the events, with $A$ the set of orderings whose consecutive pairs are arcs of $\Gamma$ (the Hamiltonian paths of $\Gamma$); when events carry categories, $\Gamma$ is the preimage of a rule set on categories. A \emph{triple grammar} forbids a set of consecutive triples. Neither has a seed or arithmetic, yet each yields a set of orderings, hence a subcomplex of the permutohedron. They are used in \cref{ch:locality} to separate what depends on arithmetic from what depends only on the \emph{locality} of the constraints. A reading with categories such as Sunny, Rainy, Snowy is a symbolic test system only; no weather data enter.

\section{Admissible sets and the permutohedral complex}\label{sec:complex}
Fix $k\ge2$, write $[k]=\{0,\dots,k-1\}$ and let $S_k$ be the set of orderings $w=(w_1,\dots,w_k)$ of $[k]$. The \emph{prefix sets} of $w$ are $I_j(w)=\{w_1,\dots,w_j\}$; the chain $I_1(w)\subset\dots\subset I_{k-1}(w)$ is the \emph{flag} of $w$, and $w\mapsto$ flag is a bijection between orderings and maximal chains of the poset $\Pm$ of proper nonempty subsets of $[k]$.

\begin{definition}[ordered partitions, faces]\label[definition]{def:faces}
A \emph{face} of the permutohedron $P_k$ is an ordered set partition $F=(B_1|\dots|B_m)$ of $[k]$ into nonempty blocks, $1\le m\le k$; it has $\dim F=k-m$. Its vertex set $V(F)$ is the set of orderings that list $B_1$ first (in any order), then $B_2$, and so on. Faces of dimension $0,1,2,3$ with block sizes $(1,\dots,1)$, $(2)$, $(2,2)$ or $(3)$, and $(2,2,2)$, $(3,2)$ or $(4)$ are vertices, edges, squares or hexagons, and cubes, hexagonal prisms or truncated octahedra. The face $F'$ is a face of $F$ if $F'$ refines $F$. The \emph{proper prefix sets} of $F$ are
\[
\Pref(F)=\{B_{<i}\cup U:\ 1\le i\le m,\ \emptyset\neq U\subseteq B_i\}\setminus\{[k]\},\qquad B_{<i}=B_1\cup\dots\cup B_{i-1}.
\]
\end{definition}
The $1$-skeleton of $P_k$ is the Cayley graph of $S_k$ for the adjacent transpositions $s_1,\dots,s_{k-1}$ (positions $j,j{+}1$); an edge is a face with one block of size two, and the graph is bipartite by the sign of the permutation. We use the standard facts about permutohedra and Coxeter groups~\cite{Ziegler,BjornerBrenti}: the faces form a regular cell decomposition of the boundary of a simple convex polytope of dimension $k-1$.

\begin{definition}[history complex]\label[definition]{def:KA}
For a set $A\subseteq S_k$ of admissible histories, the \emph{history complex} $K_A$ is the full subcomplex of $P_k$ spanned by $A$: a face $F$ of $P_k$ belongs to $K_A$ iff $V(F)\subseteq A$. Its Betti numbers $\beta_j(K_A)=\dim H_j(K_A;\F)$ are the \emph{history Betti numbers}. The 1-skeleton of $K_A$ is the induced subgraph $G_A$ of the permutohedron graph on $A$, the \emph{legal-swap graph}.
\end{definition}
For a commutative realization with \emph{distinct} labels $\ell_1<\dots<\ell_k$, labels are addressed by rank, the histories of a piece $P(m_0,Q)$ are the admissible orderings $A=A_Q(m_0)\subseteq S_k$, and $G_A$ is the swap graph of the piece, since swapping adjacent distinct labels is the transposition of adjacent positions. For contents with repeated labels the swap graph is defined but the permutohedral complex is not; Betti numbers are claimed only for distinct-label pieces.

\begin{definition}[prefix-defined admissibility]\label[definition]{def:prefix-defined}
A set $A\subseteq S_k$ is \emph{prefix-defined} if there is a family $\mathfrak F\subseteq\Pm$ such that $w\in A$ iff all proper prefix sets $I_1(w),\dots,I_{k-1}(w)$ lie in $\mathfrak F$. A prefix set $I\in\Pm$ is \emph{dead} if every ordering having $I$ as a prefix set (equivalently every flag through $I$) is not in $A$.
\end{definition}

\begin{lemma}[faces of prefix-defined complexes]\label[lemma]{lem:faces-prefix}
If $A$ is prefix-defined by $\mathfrak F$, then a face $F$ of $P_k$ lies in $K_A$ iff $\Pref(F)\subseteq\mathfrak F$.
\end{lemma}
\begin{proof}
The proper prefix sets of a vertex $w=\pi_1\cdots\pi_m\in V(F)$ ($\pi_i$ an ordering of $B_i$) are the sets $B_{<i}\cup U$ with $U$ a nonempty prefix set of $\pi_i$ (excluding $[k]$). Conversely every nonempty $U\subseteq B_i$ is a prefix set of some ordering of $B_i$ (list $U$ first), and the other blocks may be ordered arbitrarily. Hence $\bigcup_{w\in V(F)}\{I_j(w)\}=\Pref(F)$, and $V(F)\subseteq A$ iff $\Pref(F)\subseteq\mathfrak F$.
\end{proof}

\section{What requires arithmetic}\label{sec:needs}
\Cref{tab:needs} is the contract that the later chapters keep: a result in a row is claimed only for the stated class of systems. In particular the reconstruction, duality and information-ladder results hold for \emph{arbitrary} finite sets of histories, the nerve, block and cube results for every \emph{prefix-defined} system (which contains SELECT-II but is much larger), and the closure, threshold and survival laws use the divisibility lattice of SELECT-II.
\begin{table}[ht]\centering\small
\begin{tabularx}{\textwidth}{@{}>{\raggedright\arraybackslash}p{4.2cm}>{\raggedright\arraybackslash}X>{\raggedright\arraybackslash}p{2.2cm}@{}}
\toprule
Construction & Required structure & Where\\\midrule
Reverse edges, finite fibres, squarefree complete sources & SELECT-II (integers; strict growth $\ge4$ per step) & \cref{ch:ancestry}\\
Ambiguity bounds; chain-rule hierarchy $T\!\subset\!(M,Q)\!\subset\!C\!\subset\!W$ & any finite fibre; the nesting needs a \emph{commutative} realization & \cref{ch:ancestry}\\
Swap graph, accessibility profile, radius-two identity & any finite set of orderings with adjacent-transposition swaps & \cref{ch:swaps,ch:access}\\
Lemma on prefix sets (admissible iff denominators divide) & SELECT-II (divisibility); the \emph{consequence} (prefix-defined) is all that later chapters use & \cref{ch:arith}\\
Block factorization, cube-product property & any \emph{prefix-defined} $A$ & \cref{ch:arith,ch:swaps}\\
lcm closure, minimal seed, divisibility filtration, survival laws, gates & SELECT-II (divisibility lattice) & \cref{ch:arith,ch:topology,ch:atlas}\\
Reconstruction of $K_A$ from the coloured graph & arbitrary $A\subseteq S_k$ & \cref{ch:swaps}\\
Two-skeleton from the uncoloured graph & cube-product property (all prefix-defined $A$), $k\le8$ & \cref{ch:swaps}\\
Lift operator and its boundary identity & arbitrary chains on permutohedra (pure combinatorics); \emph{survival} needs arithmetic & \cref{ch:dimension}\\
Alexander duality $K_A\leftrightarrow Y_A$ & arbitrary $A\subsetneq S_k$ nonempty & \cref{ch:duality}\\
Dead-set nerve theorem & prefix-defined $A$, any $k$ for the homotopy statement; $\beta_2$ formula for $k=6$ & \cref{ch:duality}\\
Pairwise locality theorem & Hamiltonian paths of an arbitrary digraph on $k=6$ events (no arithmetic) & \cref{ch:locality}\\
\bottomrule
\end{tabularx}
\caption{Constructions and what they require. ``Prefix-defined'' includes SELECT-II (\cref{lem:F}) but is a much larger class.}\label{tab:needs}
\end{table}

\section{Probability conventions}\label{sec:probconv}
Three different probabilities occur and are never mixed. On a fibre the uniform distribution is used unless another is stated, so that $\Ent(W|T=t)=A(t)$ and every entropy of \cref{ch:ancestry} is an exact finite sum of terms $\frac nN\log_2\frac Nn$. In \cref{ch:separation} a \emph{system} (a content with a source regime) is itself a random variable and one asks how much an observation tells about a property such as $\beta_2>0$; the \emph{natural} ensemble is uniform over the systems of a stated census and the \emph{balanced} ensemble gives equal weight to the two values of the property, both being experimental conventions rather than laws of the grammar. Finally, entropies computed on a whole finite ensemble by exact class counts are \emph{in-sample}, whereas the information a rule learned on one sample gives about another is \emph{out-of-sample}; \cref{sec:predictive} shows that they differ substantially.
